% !TEX root = ../../../main.tex
\section{Integration}
\label{sec:analysis-i:integration}

\begin{pendingreview}
Integration will compare accumulated quantities with finite sums,
leading to a precise condition for their limit.
\historynote{Riemann introduced his integrability framework in his
1854 study of trigonometric representation. Dedekind published the
work posthumously in 1867~\cite[\S\S~4--6]{Riemann1867}.
Thus the integral grew partly from a question about which functions
Fourier series could represent. It broadened integration beyond
continuous functions and shaped real analysis. See \autoref{fig:portrait:riemann}.}
\historyportrait{riemann}{Bernhard Riemann}{PortraitRiemann}

\end{pendingreview}
% The definition and course proofs await the supplied lecture material.
